\documentclass[11pt]{letter}

\usepackage[margin=1in]{geometry}
\usepackage{amsmath}

\signature{Stephen Nelson\\Field Dynamics Research\\stephen@fielddynamics.org}

\address{Stephen Nelson\\Field Dynamics Research\\stephen@fielddynamics.org}

\begin{document}

\begin{letter}{The Editor\\Classical and Quantum Gravity}

\opening{Dear Editor,}

I am writing to submit the enclosed manuscript,
\textit{The Other Half of Gravity: Torsion and the
Two-Chamber Structure of the Tetrad Frame}, for
consideration as a Letter in Classical and Quantum Gravity.

This letter is a companion to the manuscript
\textit{The Coordinate-Space Geometry of the Tetrad Frame}
(CQG-115428), currently under review at your journal.
That companion letter proved that the spatial sector of the
tetrad equation produces two bounded tetrahedral regions
related by spatial parity: TetA (the identity component) and
TetB (the parity component). The present letter addresses
the question of how this two-chamber structure manifests at
the level of the affine connection.

The letter contains two formal results:

\begin{enumerate}

\item \textbf{Lemma (Section~2).} A torsion-free connection
cannot distinguish TetA from TetB. The proof follows in three
steps from the bilinearity of the metric in the frame field:
under the spatial parity that exchanges TetA and TetB, both
frame factors acquire a sign flip, leaving the metric
invariant; the Christoffel symbols, being uniquely determined
by the metric, inherit this invariance. The torsion sector is
therefore the necessary carrier of the two-chamber distinction
at the connection level.

\item \textbf{Theorem (Section~5).} Under the closure
relation $D = K - 1$, the standard torsion count
$K \times K(K{-}1)/2$ equals the product $D \times 2K$ if
and only if $K = 4$. This links the antisymmetric connection
count to the coupling range of a polynomial that organises
every standard count of the tetrad formalism.

\end{enumerate}

A supplementary note derives the TetB field fraction
$1/(x{+}1)$ from the exchange symmetry of the two chambers
combined with the rational-bijection classification on the
projective line. The dark-to-baryon mass ratio that appears
in the cold dark matter framework is shown in Section~4 to
be an algebraic identity of this partition, following from
four lines of algebra with no halo model and no free
parameters.

The manuscript does not propose a modification of general
relativity. It examines the geometric consequences of
relaxing the torsion-free specialisation within the standard
tetrad formalism. All claims are either proved locally,
derived from the companion geometry (CQG-115428), or cited
from standard references (Misner, Thorne and Wheeler 1973;
Hehl \textit{et al.}\ 1976; Cartan 1922).

Because the present letter depends directly on the
two-chamber geometry established in CQG-115428, I would be
grateful if the same referee could be considered for both
manuscripts, should that be practical. This would avoid
the need to re-evaluate the companion result within the
present review.

The manuscript has not been submitted elsewhere and is not
under consideration at any other journal.

\closing{Yours sincerely,}

\end{letter}
\end{document}
